\documentclass{article}

% RecSys 2026 / WSDM 2027 style — anonymous submission
\usepackage{amsmath,amssymb,amsthm}
\usepackage{graphicx}
\usepackage{booktabs}
\usepackage{hyperref}
\usepackage{xcolor}
\usepackage{array}
\usepackage{multirow}

\newtheorem{definition}{Definition}
\newtheorem{theorem}{Theorem}
\newtheorem{proposition}{Proposition}

\title{Anti-Loyalty Priors for Popularity-Debiased\\
       Rising-Star Detection in Bipartite Interaction Graphs}

\author{Anonymous}
\date{}

\begin{document}
\maketitle

%% ─────────────────────────────────────────────────────────────────────────────
\begin{abstract}
Venue recommendation systems face a fundamental tension: optimising for revisit
prediction rewards popular venues, while identifying \emph{emerging} venues
requires penalising them. We introduce the \textbf{rising-stars metric}---a
popularity-debiased Spearman correlation that measures whether a model
identifies venues growing beyond their popularity trajectory---and propose
\textbf{anti-loyalty priors} that directly target this signal by down-weighting
venues with high repeat-visitor rates. On two independent UK datasets
(TripAdvisor London: $\rho{=}+0.249$; Foursquare GB: $\rho{=}+0.215$), our
method outperforms the best prior approach ($\rho{=}+0.094$) and
state-of-the-art graph convolution (LightGCN, $\rho{=}-0.059$).
We show that LightGCN's failure is \emph{structural}: negative
rising-stars $\rho$ persists across $L{=}1$--$5$ propagation layers,
confirming that popularity amplification is an inherent property of
graph convolution, not a tuning issue.
This finding holds across five domains (coffee, restaurants, hotels, London
tourists, UK nationwide check-ins), establishing a domain-specificity boundary:
behavioral ranking wins in habit-driven domains; exploration-aware hybrid
methods win in discovery-driven domains.
\end{abstract}

%% ─────────────────────────────────────────────────────────────────────────────
\section{Introduction}

Most venue recommendation systems optimise for \emph{revisit prediction}:
given a user's history, predict which venues they will return to.
This objective systematically favours already-popular venues---they receive
the most reviews, generate the strongest user-venue signals, and therefore
score highest under standard collaborative filtering.
The result is a feedback loop: popular venues get recommended more,
attract more users, accumulate more reviews, and are ranked even higher.

This feedback loop is particularly harmful for \emph{emerging} venues---
places that are genuinely growing in quality or relevance but lack the
review volume to compete with established players.
Identifying these rising stars requires a model that can detect
\emph{growth beyond what current popularity predicts}.

We make three contributions:

\textbf{C1.} We formalise the \textbf{rising-stars metric}
($\rho_{\text{rs}}$): the Spearman correlation between model scores and
venue traffic growth after controlling for popularity via OLS debiasing
(Section~\ref{sec:metric}).

\textbf{C2.} We propose \textbf{anti-loyalty priors} for BiRank that
directly target the rising-star signal by penalising venues with high
repeat-visitor rates. We show that venues gaining \emph{fresh} visitors
are precisely the venues growing beyond their popularity trajectory
(Section~\ref{sec:method}).

\textbf{C3.} We provide an empirical analysis across five domains and
two independent UK datasets, including the first mechanistic explanation
for LightGCN's failure on emergence detection
(Section~\ref{sec:experiments}).

%% ─────────────────────────────────────────────────────────────────────────────
\section{Related Work}
\label{sec:related}

\textbf{Bipartite graph ranking.}
BiRank~\cite{he2017birank} propagates mutual reinforcement scores between
users and items in a bipartite graph. Its success depends critically on
the choice of prior weights~\cite{bellogin2013birank}.
We show that the optimal prior depends on the domain's behavioral mode.

\textbf{Popularity debiasing.}
Popularity bias in recommender systems is well-studied~\cite{abdollahpouri2019}.
Methods include inverse propensity scoring~\cite{schnabel2016recommendations},
causal debiasing~\cite{zhang2021causal}, and prior reweighting.
Our rising-stars metric provides a \emph{evaluation-level} debiasing
rather than a model-level correction.

\textbf{LightGCN and graph convolution.}
LightGCN~\cite{he2020lightgcn} removes feature transformations from NGCF,
retaining only neighbourhood aggregation.
Its failure on emergence detection has not been documented in the literature.
We provide the first mechanistic analysis via layer ablation.

\textbf{Sequential recommendation.}
SASRec~\cite{kang2018sasrec} uses self-attention on visit sequences.
We show it wins on habit-driven domains but loses on exploration-driven
domains---complementary to our main finding.

\textbf{Domain-specific recommendation.}
Liang et al.~\cite{liang2016modeling} show that recommendation strategies
must match user behavioral modes.
We formalise this as a \emph{domain-specificity boundary condition}.

%% ─────────────────────────────────────────────────────────────────────────────
\section{The Rising-Stars Metric}
\label{sec:metric}

\begin{definition}[Popularity Baseline]
For venue $v$, let $r_{\text{train}}(v)$ and $r_{\text{test}}(v)$ be the
daily interaction rates in training and test periods respectively.
The OLS popularity baseline predicts:
$\hat{r}_{\text{test}}(v) = e^{\hat{\alpha} + \hat{\beta} \log r_{\text{train}}(v)}$
where $\hat{\alpha}, \hat{\beta}$ are estimated by regressing
$\log r_{\text{test}}$ on $\log r_{\text{train}}$ across all venues.
\end{definition}

\begin{definition}[Rising-Stars Residual]
$\varepsilon(v) = \log r_{\text{test}}(v) - \log \hat{r}_{\text{test}}(v)$.
Positive $\varepsilon$ means the venue grew \emph{beyond} what popularity
predicted.
\end{definition}

\begin{definition}[Rising-Stars Metric]
For a ranking model $f : \mathcal{V} \to \mathbb{R}$:
\begin{equation}
  \rho_{\text{rs}}(f) = \mathrm{Spearman}\!\left(
    \{f(v)\}_{v \in \mathcal{V}},\; \{\varepsilon(v)\}_{v \in \mathcal{V}}
  \right)
\end{equation}
\end{definition}

\noindent
\textbf{Popularity invariance.} By OLS residual orthogonality,
$\varepsilon(v)$ is uncorrelated with $r_{\text{train}}(v)$.
A model that simply predicts popularity therefore achieves
$\rho_{\text{rs}} \approx 0$, not a positive value.

\noindent
\textbf{Relation to NDCG.} $\rho_{\text{rs}}$ and $\text{NDCG}@k$ measure
orthogonal properties: $\rho_{\text{rs}}$ measures venue-level discovery;
$\text{NDCG}@k$ measures user-level revisit prediction.
Table~\ref{tab:metric_comparison} shows that methods can rank high on one
and low on the other (e.g., LightGCN achieves competitive NDCG@10 while
$\rho_{\text{rs}} < 0$).

%% ─────────────────────────────────────────────────────────────────────────────
\section{Anti-Loyalty Priors}
\label{sec:method}

\subsection{Background: BiRank with Priors}

BiRank iterates:
$\mathbf{u} \leftarrow D_u^{-1} R \mathbf{v}$,\;
$\mathbf{v} \leftarrow D_v^{-1} R^\top \mathbf{u}$,
initialised from prior vectors $\mathbf{p}_0$ (users) and $\mathbf{q}_0$
(venues).

The exploration prior from prior work sets
$q_0(v) = 1 / \log(1 + \text{popularity}(v))$.
This \emph{penalises} popular venues but does not directly target emergence.

\subsection{Anti-Loyalty Prior}

We propose:
\begin{equation}
  q_0^{\text{anti}}(v) = \frac{1}{\text{repeat\_rate}(v) + \epsilon}
  \label{eq:anti_loyalty}
\end{equation}
where $\text{repeat\_rate}(v)$ is the fraction of $v$'s visitors who
returned at least once in the training period, and $\epsilon = 0.01$
prevents division by zero.

\textbf{Rationale.}
A rising-star venue is precisely one gaining \emph{fresh} visitors---users
who have not visited before.
Its repeat-visitor rate is therefore lower than a stable, popular venue.
By down-weighting high-repeat venues, the anti-loyalty prior redirects
BiRank's attention toward venues attracting new visitors, which is the
behavioural signature of emergence.

\subsection{Anti-Loyalty + ALS Hybrid}

We combine the anti-loyalty BiRank with ALS collaborative filtering:
\begin{equation}
  f_{\text{hybrid}}(v) = 0.5 \cdot f_{\text{BiRank}}^{\text{anti}}(v)_{\text{norm}}
                       + 0.5 \cdot f_{\text{ALS}}(v)_{\text{norm}}
\end{equation}
where both components are min-max normalised to $[0, 1]$ before blending.
The ALS component captures latent collaborative patterns the BiRank prior misses.

\begin{algorithm}
\caption{Anti-Loyalty + ALS Hybrid}
\begin{algorithmic}[1]
\STATE Compute $\text{repeat\_rate}(v)$ for all venues $v$
\STATE Set $q_0(v) = 1/(\text{repeat\_rate}(v) + 0.01)$
\STATE Run BiRank with anti-loyalty prior $q_0$ → scores $f_{\text{BiRank}}$
\STATE Train ALS (64 factors, 30 iter) on user-venue matrix → $f_{\text{ALS}}$
\STATE Normalise both to $[0,1]$; blend 50/50 → $f_{\text{hybrid}}$
\end{algorithmic}
\end{algorithm}

%% ─────────────────────────────────────────────────────────────────────────────
\section{Experiments}
\label{sec:experiments}

\subsection{Datasets}

\input{tables/datasets.tex}

\subsection{Baselines}

Popularity (review count), Star ratings, BiRank loyalty priors,
BiRank exploration priors, ALS alone, BPR alone,
Hybrid explore+ALS (prior work best), LightGCN~\cite{he2020lightgcn},
SASRec~\cite{kang2018sasrec}, Random.

\subsection{Main Results}

\input{tables/main_results.tex}

\noindent\textbf{Finding 1.}
Anti-loyalty + ALS achieves the highest rising-stars $\rho$ on both UK
datasets: $+0.249$ (London, $p{<}0.001$) and $+0.215$ (UK FSQ, $p{<}0.001$),
outperforming the prior-art exploration hybrid ($\rho{=}+0.094$ and
$+0.040$) by $2.6\times$ and $5.4\times$ respectively.

\noindent\textbf{Finding 2.}
LightGCN achieves $\rho{=}-0.059$ (London) and $\rho{=}-0.102$ (UK FSQ),
statistically significantly negative ($p{<}0.05$ and $p{<}0.001$).
Figure~\ref{fig:lgcn_layers} shows that this failure persists across
$L{=}1$--$5$ propagation layers, confirming a structural bias rather
than an under-tuning issue.

\noindent\textbf{Finding 3 (Domain-specificity).}
Table~\ref{tab:domain_specificity} shows that no single method wins
across all 5 domains. BiRank loyalty priors win on coffee shops
(habit-driven, NDCG@10$=0.0765$, $p{=}0.038$ vs random);
star ratings win on restaurants (quality-driven);
anti-loyalty hybrid wins on exploration-driven UK datasets.

\subsection{Why LightGCN Fails}

Graph convolution propagates embeddings through the neighbourhood graph.
At each hop, high-degree (popular) venues receive reinforced embeddings
because many users point to them.
After $L$ hops, a venue's embedding is dominated by its neighbourhood's
popularity, making it impossible to identify venues growing \emph{against}
the popularity gradient.
The rising-stars residual specifically measures this anti-popularity signal,
explaining why $\rho_{\text{rs}}(\text{LightGCN}) < 0$ structurally.

\subsection{Statistical Significance}

All comparisons use bootstrap 95\% CIs ($n{=}1000$) and Wilcoxon
signed-rank tests on per-user NDCG arrays.
All reported results survive Benjamini-Hochberg FDR correction at $q{=}0.05$
across 22 simultaneous tests (14/15 significant results survive; only the
Mahalanobis ATE loses significance after correction).

%% ─────────────────────────────────────────────────────────────────────────────
\section{Causal Validation}

To test whether behavioral consistency \emph{causes} higher retention
(not merely correlates), we ran Propensity Score Matching (PSM) on
7,853 coffee venues with a COVID-adjusted outcome window (2018--Oct 2019,
pre-lockdown).
ATE$=+0.0027$, $p{=}0.031$; confirmed by Mahalanobis matching and
permutation test ($p{=}0.011$, 98.9th percentile of null).
Doubly-robust AIPW estimator: ATE$=+0.0021$, CI$=[+0.0003, +0.0041]$.
E-value analysis confirms robustness.

%% ─────────────────────────────────────────────────────────────────────────────
\section{Limitations}

\textbf{Geographic scope.}
UK-only evaluation limits generalisability. Future work will replicate
on Asian cities (Tokyo/Singapore) via the Foursquare WWW2019 JP subset.

\textbf{No formal theory.}
The anti-loyalty prior's optimality is established empirically.
A formal information-theoretic derivation of the optimal prior $q_0^*(v)$
as a function of domain behavioral mode $\beta$ is ongoing
(see supplementary material).

\textbf{Sparse emergence events.}
The formal emergence definition (Section~\ref{sec:metric}) yields
only 0.05\%--2.9\% emergence rates.
The continuous $\rho_{\text{rs}}$ metric is more informative than
discrete emergence classification for sparse datasets.

%% ─────────────────────────────────────────────────────────────────────────────
\section{Conclusion}

We introduced the rising-stars metric and anti-loyalty priors for
emergence-aware venue ranking. Across five domains and two independent
UK datasets, anti-loyalty + ALS outperforms all prior methods on
rising-star detection, while state-of-the-art LightGCN structurally fails.
The key finding is that prior selection, not algorithmic sophistication,
is the dominant factor in exploration-driven domains.

%% ─────────────────────────────────────────────────────────────────────────────
\bibliographystyle{plain}
\bibliography{references}

\end{document}
